% !PRINTED 142
\weberpage{142}
%<<CONT
den beiden Invarianten $p, p$ oder mit einer Invariante $p^2$), und
wenn $p$ von $q$ verschieden ist, eine Gruppe (mit den beiden Invarianten
 $p, q$) giebt.

Es bleibt uns also noch der Fall
\[
p \equiv 1\ (\mathrm{mod}\ q)
\]
zu untersuchen, in dem $\nu$ jeden der Bedingung
\[
\nu \equiv 0\ \left(\mathrm{mod}\ \frac{p-1}{q}\right)
\]
genügenden Werth haben kann, deren es $q$ nach dem Modul
$p - 1$ verschiedene giebt, nämlich:
\weq{7}{\nu = 0,\quad \frac{p-1}{q},\quad 2\,\frac{p-1}{q},\quad \dots\ \frac{(q-1)(p-1)}{q}.}
Der Fall $\nu = 0$ führt auch hier auf eine \gesperrt{Abel'sche} Gruppe.
Es sind aber auch die anderen Werthe von $\nu$ zulässig, die zu
ebenso vielen nicht commutativen Gruppen führen. Diese nicht
commutativen Gruppen sind unter einander isomorph und werden
auf einander zurückgeführt, wenn man $b$ durch ein $b^\beta$ ersetzt,
wie die Formel (5) zeigt.

Wir können $b^q = 1$ annehmen, denn es ist gewiss $b^q$ in $Q$
enthalten und $b^{pq} = 1$. Wenn also $b^q$ nicht $= 1$ ist, so ersetzen
wir $b$ durch $b^p$.

Man erhält eine der nicht commutativen Gruppen, wenn
man in
\weq{8}{\Theta = a^\alpha b^\beta}
$\alpha$ die Reihe der Zahlen $0, 1, \dots p - 1$ und $\beta$ die Zahlen
$0, 1, \dots q - 1$ durchlaufen lässt. Die Zusammensetzung zweier
Elemente dieser Gruppe ergiebt sich nach (5) aus
\weq{9}{b^\beta a^\alpha = a^{\alpha g^{-\nu\beta}} b^\beta, \quad a^p = 1,\ b^q = 1,}
wodurch man jedes Compositum aus Elementen $\Theta$ auf die Form
$\Theta$ zurückführen kann.

Um zu zeigen, dass die Elemente (8) nach den Compositionsregeln
 (9) wirklich eine Gruppe bilden, hat man die Eigenschaften
der Gruppe, nämlich, dass aus $\Theta\Theta' = \Theta\Theta''$ und aus $\Theta'\Theta = \Theta''\Theta$
folgt, dass $\Theta' = \Theta''$ ist, und ferner das associative Gesetz
\[
\Theta(\Theta'\Theta'') = (\Theta\Theta')\Theta''
\]
nachzuweisen. Beides aber ergiebt sich sehr leicht aus der Zusammensetzung,
die aus (9) folgt:
\weq{10}{a^\alpha b^\beta a^{\alpha'} b^{\beta'} = a^{\alpha+\alpha' g^{-\nu\beta}} b^{\beta+\beta'}.}
